\clearpage
\subsection{Sudamérica}

\subsubsection{MERCOSUR}

El \gls{mercosur} es el proceso de integración comercial más importante de América Latina \parencite[p.~195]{Lakos2025}. Fue creado inicialmente por Argentina, Brasil, Paraguay y Uruguay, con la Tratado de Asunción en 1991. En 2006 Venezuela también se adhirió al tratado constitutivo, aunque desde 2017 se encuentra suspendida en todos sus derechos y obligaciones de forma indefinida \parencite[p.~769]{EuropaDirectory2022}. Posteriormente, Bolivia se adhirió al Tratado de Asunción en 2015, aunque tardó 8 años en finalizar el proceso de adhesión \parencite[p.~199]{Lakos2025}. 
%Posteriormente, también lo hizo Bolivia en 2015.
El \gls{mercosur} también considera una membresía como Estados asociados, la cuál esta orientada a los países de la \gls{aladi}\footnote{La \gls{aladi} es un organismos regional creado en 1980 integrada por 13 países: Argentina, Bolivia, Brasil, Chile, Colombia, Cuba, Ecuador, México, Panamá, Paraguay, Perú, Urugua y Venezuela, que tiene como objetivo final la creación de un mercado común latinoamericano.} con los cuales el MERCOSUR tiene acuerdos de libre comercio. Actualmente, los países con este tipo de membresía son Chile, Colombia, Ecuador, Guyana, Panamá, Perú y Surinam \parencite{MERCOSUR_Paises}.

% La estructura institucional sugiere una fuerte inspiración de las Comunidades Europeas en ese momento, pero el MERCOSUR, a diferencia de la integración europea, no ha pasado de un enfoque principalmente intergubernamental a una estructura de toma de decisiones más federal que no siempre requiere el acuerdo de todos los Estados miembros en todo. Lako 198

La zona de libre comercio del \gls{mercosur} entró en vigor a partir de inicios de 1995, con la eliminación de los aranceles en el 85\% del comercio intrarregional. Además, para un conjunto de productos especiales se acordó un régimen de eliminación gradual de aranceles. Ese mismo año, también entró en vigor la unión aduanera del \gls{mercosur}, con un arancel común del 0\% al 20\%. En 1996 los países miembros firmaron un protocolo separado para la liberalización de los servicios y las compras gubernamentales durante un periodo de 10 años. Durante la cumbre del \gls{mercosur} de 2003, los jefes de Estado de los cuatro países miembros acordaron fortalecer la integración del bloque regional y armonizar los aranceles de importación para 2006, creando así las bases para un mercado único \parencite[p.~770]{EuropaDirectory2022}. En 2017, los Estados miembros firmaron un protocolo para la cooperación y la facilitación de inversiones entre los miembros. Además de la liberalización económica y el establecimiento de la unión aduanera, el \gls{mercosur} ha reducido significativamente las barreras al movimiento de la fuerza laboral dentro de la región. Los residentes de los Estados miembros pueden vivir y trabajar en cualquier otro país miembro. Pese a ello, esta organización no ha logrado convertirse en realmente en un mercado común, en opinión de \textcite{Lakos2025}. Una de las razones que ha contribuido a que la integración del \gls{mercosur} no se profundice son las barreras comerciales no arancelarias, como los cargos comerciales altos debido a una infraestructura poco desarrollada y las presiones regulatorias sobre los servicios comerciales \parencite[p.~201]{Lakos2025}.

La estructura del comercio internacional del \gls{mercosur} se caracteriza por un predominio de exportaciones de productos primarios al exterior del bloque; en tanto que el comercio intrarregional es más diverso e incluye una gran proporción de bienes manufacturados \parencite[p.~201]{Lakos2025}. 

%or otro lado, sin embargo, el Mercado Común del Sur no ha podido expandirse significativamente ni en anchura mediante la ampliación de su número total de miembros –ni en profundidad– yendo más allá de su carácter estrictamente intergubernamental. p. 2000

Desde 1995, Bolivia y Chile se integraron como miembros asociados. En el caso de Bolivia, desde 1997 forma parte de la zona libre comercio del \gls{mercosur}, pero no de su unión aduanera. En 1996 Chine firmó un Acuerdo de Complementación Económica por el que se eliminaban los aranceles sobre la mayoría de los productos en un plazo de 10 años. Aunque al igual que Bolivia, Chile también quedaba fuera de la unión aduanera, los acuerdos establecían su participación en proyectos de infraestructura diseñados para dar a los países del \gls{mercosur} acceso al Océano Pacífico. Sin embargo, el procedimiento para el ingreso de Chile como miembro de pleno derecho del \gls{mercosur} se suspendieron en el año 2000, después del anuncio de conversaciones bilaterales entre Chile y Estados Unidos para la firma de un tratado de libre comercio. Perú se conviritó en miembro asociado del MERCOSUR en 2003; Colombia y Ecuador hicieron lo mismo al año siguiente \parencite[p.~770]{EuropaDirectory2022}.

El \gls{mercosur} también ha establecido acuerdos de libre comercio multileratelaes. En 2005 entró un acuerdo de libre comercio con la Comunidad Andina que preveía la eliminación gradual de los aranceles sobre el comercio de los miembros de ambos países en un plazo de 14 años. Además, ambos organismos acordaron extender la membresía asociada a todos los Estados miembros de ambos grupos \parencite[p.~771]{EuropaDirectory2022}. Recientemente, en enero de 2026 se firmó un acuerdo de asociación entre la \gls{ue} y el \gls{mercosur}. Mediante este acuerdo, los países del \gls{mercosur} obtienen acceso preferencial al mercado europeo, con la eliminación de aranceles para el 92\% de las exportaciones \parencite{MERCOSUR_UE}.

En sintesis, bajo la escala de integración de Balassa, el \gls{mercosur} puede considerarse como una zona de libre comercio que garantiza la libre circulación de bienes, servicios, capital y fuerza de trabajo, así como una unión aduanera en pleno funcionamiento. Sin embargo, no ha alcanzado el nivel de una unión económica y monetaria, ni de un mercado único \parencite[p.~206]{Lakos2025}. 


\subsubsection{Comunidad Andina (CAN)}

La \gls{can} fue fundada en 1969 por cinco países latinoamericanos: Bolivia, Chile, Colombia, Ecuador y Perú, bajo el nombre de Pacto Andino. Chile se retiró de la Comunidad Andina en 1976, bajo la dictadura militar de Augusto Pinochet. Venezuela también salio de la comunidad en 2006 para ingresar al \gls{mercosur} \parencite[pp.~223]{Sanchez2025}. En 1993 se creó la Zona Andina de Libre Comercio entre Bolivia, Colombia, Ecuador y Venezuela; Perú se incorporó a finales de 2005 y Venezuela continuó participando hasta 2011, pese a su salida de la Comunidad Andina \parencite{CAN_50anos}. En 1996 se fortalece institucionalmente la integración andina mediante el Acuerdo de Cartagena, que entre otros elementos, decide renombrar el proceso de integración como Comunidad Andina. En 2001, se avanza hacia un proceso de libre tránsito, permitiendo viajar por la región con documentos de identidad nacional, sin pasaporte ni visa. En 2017 se lanza un satelite de la Comunidad Andina y se aprueba un marco regulatorio de mercados de energía eléctrica, mostrando avances en la integración de la infraestructura regional \parencite{CAN}. 

Un elemento fundamental en la configuración institucional de la Comunidad Andina fue el establecimiento del Sistema Andino de Integración en 1996, conocido como Protocolo de Trujillo. Uno de los puntos importantes de las reformas derivadas de este protocolo fue la adopción del principio de supranacionalismo, que permite que las leyes andinas sean vinculantes para los Estados miembros, y no requieran de procesos especiales para entrar en vigor \parencite[pp.~224--225]{Sanchez2025}.

Actualmente, la Comunidad Andina es una zona de libre comercio consolidada, con la aspiración de formar un mercado común, previsto inicialmente para el año 2020 \parencite[p.226]{Sanchez2025}. Sin embargo, presenta niveles relativamente bajos de comercio intrarregional, en gran parte debido a que los cuatro países exportan principalmente bienes primarios o sus derivados. Como ejemplo, las principales destinos de las exportaciones de los cuatro miembros son Estados Unidos, China y la Unión Europeo, acumulando el 49.3\% de las exportaciones totales de la Comunidad Andina. En contraste, la exportaciones intrarregionales solo representan el 6.4\% de las exportaciones de los países de la Comunidad Andina \parencite[p.227]{Sanchez2025}. 


\subsubsection{Alianza del Pacífico}

La alianza del Pacífico es otra iniciativa de integración regional latinoamericana, formada por Colombia, Chile, México y Perú desde 2011. Su antecedente directo es el Foro del Arco del pacífico Latinoamericano impulsado por el gobierno peruano desde 2006 con el objetivo de fortalecer el comercio con Asia. Dentro de sus objetivos originales se encuentra aumentar la integración económica y avanzar gradualmente hacia la libre circulación de bienes, servicios y productos en la región. A partir de 2014, se eliminaron más del 95\% de los aranceles entre los páises miembros. Además, los cuatro miembros tienen \glspl{tlc} con Canadá, la Unión Europea y Estados Unidos, así como con varios países asiáticos. Además, la Alianza ha firmado acuerdos comerciales con China e India \parencite[pp.~250--251]{Ertl2025}. 

Al igual que la \gls{can}, existe un bajo nivel relativo de comercio interno. El porcentaje que representan las exportaciones hacia otro país de la Alianza respecto al total de exportaciones nacionales va desde el 1.38\% en el caso de México, hasta el 8.15\% de Colombia, con datos del 2021. El bajo comercio intrarregional y la estrategia \textit{hacia afuera} por el caracter estructural de sus exportaciones. En el caso de Chile, Colombia y Perú, los productos de exportación son principalmente materias primas; en cambio, México tiene un comercio más diversificado y dirigido hacia el mercado estadounidense \parencite[p.~252]{Ertl2025}. Sin embargo, el comercio intrarregional es mas importante en términos de servicios. Las exportaciones de servicios a otros países de la Alianza representa el 28.3\% de las exportaciones de servicios de Perú. En Chile y Colombia, esta proporción asciende al 10.1\% y el 17.8\% respectivamente. México es el único país de la alianza que presenta un bajo nivel en de exportaciones en servicios intrarregionales, equivalente al 4.5\% del total de sus exportaciones de servicios. Los principales servicios comercializados dentro de la Alianza son el transporte y los viajes \parencite[p.~253]{Ertl2025}.

Un elemento importante dentro de la alianza es la libre circulación de personas. En 2012, México abolió las visas para los ciudadanos colombianos, peruanos y chilenos. Perú hizo algo equivalente al año siguiente, aunque con algunas restricciones. De esta forma, desde 2013 se abolieron las frotneras técnicas a la libre circulación entre los cuatro países miembros \parencite[p.~253]{Ertl2025}.

Em opinión de \textcite{Ertl2025}, la Alianza del Pacífico sigue una estrategia \textit{orientada hacia el exterior}, principalmente hacia la región de Asia y el Pacífico. %Como tal, podemos distinguir una parte de la Alianza donde el comercio con Estados Unidos y Europa es más significativo, mientras que Chile y Perú, China y el mercado asiático son las fuentes dominantes de ingresos comerciales. p.256-257Sin embargo, en la región de Asia y el Pacífico, entre todos los miembros, se puede afirmar que China es el socio comercial más importante,

%The Pacific Alliance, also treated in this book (see Chapter 10), is one of the most recent integration platforms, combining trade liberalisation between its members with initiatives such as visa-free tourist travel, a common stock exchange, and joint embassies (Economist, 2015) p. 197 



%There are more than 15 platforms of regional cooperation in Latin America (Ayuso & Villar, 2014). Most countries take part in more than one organisation. These institutions sometimes overlap in their different functions (economic, political, social, etc.). That is why this latest wave is often referred to as interregionalism. The current picture of Latin America’s integration indicates a proliferation and oversupply of regional institutions (Van Klaveren, 2017) (Figure 8.1). p. 197 

\subsubsection{Análisis de comunidades}



\begin{figure}[H]
	\centering
	\includegraphics[width=1\textwidth]{sankey_diagrams/south_america.png}
	\caption{Evolución de la pertenencia a comunidades de los sectores económicos de países en Sudamérica: Brasil, Argentina, Chile, Colombia y Perú.}
	\label{DiagramaSankey_Sudamérica}
\end{figure}

%\begin{figure}[H]
%	\centering
%	\begin{subfigure}{.9\textwidth}
	%		\centering
	%		\includegraphics[width=\textwidth]{images/maps/america_1995_1.png}
	%		\caption{1995}
	%	\end{subfigure}
%	\begin{subfigure}{.9\textwidth}
	%		\centering
	%		\includegraphics[width=\textwidth]{images/maps/america_2008_1.png}
	%		\caption{2008}
	%	\end{subfigure}
%	\begin{subfigure}{.9\textwidth}
	%		\centering
	%		\includegraphics[width=\textwidth]{images/maps/america_2008_1.png}
	%		\caption{2022}
	%	\end{subfigure}
%	\caption{Mapa de las comunidades identificadas en Sudamérica$^{*}$.}
%	\label{Mapa_Comunidades_Sudamerica}

%	\caption*{\footnotesize $^{*}$ Los sectores económicos se disponen circularmente en torno al centro de cada país al que pertenecen, y se ordenan en el sentido de las manecillas del reloj. El color de cada sector identifica la comunidad a la que pertenece, denominada según el país que la encabeza. Las aristas representan únicamente los flujos comerciales cuya magnitud supera el percentil 50 de la distribución regional, y su grosor es proporcional a la intensidad del flujo correspondiente.}
%\end{figure}
